\section{17.1 언어모델링: 다음 토큰 예측과 사전학습·파인튜닝}
\begin{frame}{언어모델: ``다음 토큰을 예측하는 함수''}
  \begin{itemize}
    \item LLM은 텍스트를 \kb{토큰}(token)의 시퀀스로 다룬다 —
      단어 전체이거나, \textbf{더 흔하게} 단어의 일부 조각.
    \item 사전학습의 목표: 지금까지의 토큰이 주어졌을 때,
      다음 토큰의 확률분포를 예측 ---
      \[P(x_t \mid x_1, x_2, \ldots, x_{t-1})\]
    \item ``프랑스의 수도는 \_\_\_''를 잘 완성하려면 모델은
      ``파리''라는 사실을 어딘가에 담아둬야 한다.
    \item 이런 문장이 수십억 개 쌓이면, 단순한 ``다음 단어
      맞히기''를 잘 하려는 모델이 \textbf{문법·사실 지식·추론
      능력}까지 자연스럽게 습득.
    \item 확률은 \kb{decoder-only Transformer}로 계산 — 미래
      토큰을 보지 못하게 \textbf{마스킹}한 Week 12의 트랜스포머.
      ``표현하는 장치'' $\to$ ``확률분포를 계산하는 실제 엔진''.
    \item 학습: 실제 다음 토큰과 예측 분포 사이의
      \kb{교차 엔트로피 손실}(Week 2.3 정보이론에서 나온 것과
      같은 형태) 최소화.
    \item 토큰 단위 평균 $-\log p$에 지수 $e$를 취하면
      \kb{퍼플렉시티}($\text{PPL} = e^{\mathcal{L}}$) — ``다음
      토큰을 고를 때 평균 몇 개의 선택지를 두고 헤매는가''.
      낮을수록 확신 있는 예측.
  \end{itemize}
\end{frame}
